\documentclass[12pt,a4paper]{article}
\usepackage[spanish]{babel}
\usepackage[utf8]{inputenc}
\usepackage[T1]{fontenc}
\usepackage{amsmath, amssymb}
\usepackage{enumitem}
\usepackage{geometry}
\geometry{margin=2.5cm}

\title{Prueba Individual de Modulo de Pensamiento Sistémico\\
\large Modelos Y Arquitectura de Sistemas.}
\author{respuestas personales, no lo tome como si fueran ciertas - Pitehr Hurtado-Cayo}

\begin{document}

\maketitle

\begin{enumerate}[label=\textbf{Pregunta \arabic*.}]
    \item \textbf{Comente. "Sinergia es equivalente a interaccion entre las partes".}

          \textbf{Respuesta:} \\
          FALSO. La sinergia se refiere a una propiedad emergente, es decir surge/distingue a partir de las interacciones de las partes de un sistema,
          Por ejemplo: si conceptualizamos un sistema como la sala de clases, podemos distinguir dos subsistemas, alumno y profesor, ambos interactuan entre si (charla, debate, exposiciones, etc)
          pero la sinergia se refiere a la propiedad emergente de aprendizaje, que surge a partir de las interacciones entre alumno y profesor, pero no es equivalente a la interacción entre las partes.

    \item \textbf{Explique por que estabilidad/inestabilidad de un sistema es un fenomeno que se puede explicar sobre la existencia de relaciones causales.}

          \textbf{Respuesta:} \\
          La estabilidad o inestabilidad de un sistema se puede explicar a través de las relaciones causales, dado que
          un sistema es estable o inestable dependiendo de como reacciona el sistema frente a perturbaciones del entorno,
          es decir, si en el entorno sucede A el sistema se mueve al estado B, entonces dependiendo del estado B se puede clasificar al sistema como estable o inestable.
          Si no existiera reaccion frente a cambios del entorno, entonces estamos frente a un sistema indiferente, es decir, el estado del sistema no se ve afectado por cambios en el entorno, lo cual es un caso particular de estabilidad, pero no es la única forma de estabilidad, dado que un sistema puede ser estable a pesar de reaccionar frente a cambios del entorno, siempre y cuando el estado del sistema se mantenga dentro de ciertos límites.

    \item \textbf{¿Es el comportamiento teleologico una propiedad sinergica del sistema?}

          \textbf{Respuesta:} \\
          Si, el comportamiento teleológico es una propiedad sinérgica del sistema, dado que el comportamiento teleológico se refiere a la capacidad de un sistema para alcanzar un objetivo o propósito específico, y esta capacidad surge a partir de las interacciones entre las partes del sistema, es decir, no es una propiedad de una sola parte del sistema, sino que es una propiedad emergente que surge a partir de la interacción de todas las partes del sistema.

    \item \textbf{Comente: "En la base del modelo del proceso operacional de Acevedo, se encuentra un sistema de causalidad circular".}

          \textbf{Respuesta:} \\
          El modelo de hector Acevedo no establece causalidad circular, dado que este hacer referencia a un anillo cerrado, donde existe retroalimentacion.
          Es decir, dado eventos, el conductor establece ajustes al sistema, no establece causalidad entre eventos y ajustes/vector de control.

    \item \textbf{¿Tiene sentido referirse a la Ley de la Variedad Requerida de Ashby en la operacion de un sistema homeostatico?. Explique.}

          \textbf{Respuesta:} \\
          Un sistema homeostático es aquel que mantiene su estabilidad a través de mecanismos de retroalimentación, es decir, el sistema reacciona frente a cambios del entorno para mantener su estado dentro de ciertos límites.
          La Ley de la Variedad Requerida de Ashby establece que para que un sistema pueda mantener su estabilidad, debe tener una variedad de respuestas que sea al menos igual a la variedad de perturbaciones que puede enfrentar el sistema, es decir, el sistema debe ser capaz de reaccionar frente a cualquier perturbación que pueda enfrentar, para mantener su estabilidad.
          Por lo tanto, tiene sentido referirse a la Ley de la Variedad Requerida de Ashby en la operación de un sistema homeostático, dado que esta ley establece un principio fundamental para el funcionamiento de sistemas homeostáticos.
          Un ejemplo: sistema "vehiculo en movimiento" donde distinguimos dos subsistemas "conductor" y "vehiculo" entonces el conductor debe tener al menos la variedad de respuestas que el vehiculo puede enfrentar, es decir, el conductor debe ser capaz de reaccionar frente a cualquier perturbación que pueda enfrentar el vehiculo, para mantener la estabilidad del sistema "vehiculo en movimiento".


\end{enumerate}

\end{document}